%! Setting Up a Test for a Population Mean — Unit 7, Lesson 7.4 — Homework (Answer Key)
\documentclass[10pt]{article}
\usepackage{apstats-article}
\usepackage{apstats-key}   % pulls in -boxes; do NOT also load -boxes
\newcommand{\TallMath}[1]{$\displaystyle #1\rule[-1.4em]{0pt}{3.2em}$}

\begin{document}

\pageheader{Unit 7, Lesson 7.4}{Homework --- Answer Key}
\namedateperiod

\vspace{0.15in}

\textit{For each problem: (a) name the inference method and df; (b) state $H_0$ and $H_a$;
(c) verify both conditions. Show all work.}

\begin{enumerate}

\item \textbf{Water Temperature.}

\begin{enumerate}[label=\alph*.]
  \item Method: \ans{one-sample $t$-test for $\mu$} \quad $df = $ \ans{27}
  \item $H_0$: \ans{$\mu = 77$} \quad $H_a$: \ans{$\mu > 77$}
  \item Conditions:
        \begin{itemize}
          \item Random: \ansline{28 readings were randomly collected --- met.}
          \item Normal ($n < 30$): \ansline{$n = 28 < 30$; data are slightly right-skewed with no outliers. Slight skewness with no outliers is generally acceptable for small samples --- Normal condition met.}
        \end{itemize}
\end{enumerate}

\item \textbf{Delivery Times.}

\begin{enumerate}[label=\alph*.]
  \item Method: \ans{one-sample $t$-test for $\mu$} \quad $df = $ \ans{54}
  \item $H_0$: \ans{$\mu = 3$} \quad $H_a$: \ans{$\mu \ne 3$}
  \item Conditions:
        \begin{itemize}
          \item Random: \ansline{Random audit of 55 orders --- met.}
          \item Normal: \ansline{$n = 55 \ge 30$; CLT applies regardless of shape --- met.}
        \end{itemize}
\end{enumerate}

\item \textbf{Matched Pairs: Reaction Time.}

\begin{enumerate}[label=\alph*.]
  \item $d_i = $ \ans{before} $-$ \ans{after} \quad; \quad $\mu_d = $ \ans{true mean difference in reaction time (before $-$ after) for all participants using the app}
  \item Method: \ans{matched-pairs $t$-test for $\mu_d$} \quad $df = $ \ans{9}
  \item $H_0$: \ans{$\mu_d = 0$} \quad $H_a$: \ans{$\mu_d > 0$}
  \item Conditions:
        \begin{itemize}
          \item Random/Independence: \ansline{10 participants; assume a random assignment or representative sample. Differences are independent if participants are independent.}
          \item Normal ($n < 30$): \ansline{$n = 10 < 30$; distribution of differences is roughly symmetric with no outliers --- met.}
        \end{itemize}
\end{enumerate}

\item \textbf{AP-Style.}

\begin{enumerate}[label=(\Alph*)]
  \item The condition is met because $n = 12$ is large enough for the CLT.
  \item The condition is not met; moderate right skewness with $n < 30$ violates the condition.
  \item The condition is met; moderate skewness with no outliers is acceptable for $n < 30$. \textcolor{keyred}{\textbf{$\leftarrow$ correct}}
  \item The condition cannot be evaluated without knowing the population standard deviation.
\end{enumerate}
Answer: \ans{C} \quad Brief justification: \ansline{When $n < 30$, we check that the sample shows no strong skewness or outliers. Moderate skewness without outliers is generally acceptable. (A) is wrong because $n = 12$ is not large enough for CLT.}

\end{enumerate}

\begin{extensionbox}
Extension answers will vary. Key ideas: nonparametric tests (sign test, Wilcoxon signed-rank)
make fewer assumptions about the population distribution; trade-off is that they are generally
less powerful (less likely to detect a true effect) than the $t$-test when the $t$-test's
assumptions are satisfied.
\end{extensionbox}

\end{document}
